% 部署与复现扩展（中文）
\subsection{部署清单与操作流程}
\label{sec:deploy-checklist-zh}

\paragraph{诊前标定.}
(1) 通过 \texttt{import\_shezhenv3.py} 导入 ShezhenV3 或本地划分；
(2) 启用 ROI 运行 \texttt{calibrate\_tau\_tcm.py}，将 \texttt{recommended\_tau.json} 部署至边缘网关；
(3) 可选：在具备 GPU/NN 运行时训练 B3；
(4) 合并日志前用 \texttt{validate\_jsonl.py} 校验 JSONL 模式。

\paragraph{运行时参数.}
环境变量 \texttt{PAPER1\_VISION\_THRESHOLD\_TAU} 可覆盖 JSON 中的 $\tau$；\texttt{use\_roi} 默认 true；重采预算 $K{=}2$；Edge-IQA 固定 $512$ 缩放以保证延迟稳定。

\paragraph{操作员可读标志.}
\texttt{blur} 提示患者保持静止并重新张口；\texttt{underexposed}/\texttt{overexposed} 建议调整环光后再调用 FR3 技能。
日志保留 \texttt{q\_img}、路由决策与重试次数，拒帧不必上传云端。

\paragraph{硬件分级.}
\textbf{A 级（纯 CPU 网关）：}B2 + LangGraph，IQA p95 $\sim$10\,ms。
\textbf{B 级（GPU 边缘）：}可选 B3，以可解释性换取更强模糊排序。
\textbf{C 级（Franka 辅轨）：}M6 RTT，不与表 II 离线矩阵直接数值对比。

\paragraph{多中心部署.}
各中心应在本地验证批（$\geq 200$ 清晰/模糊对）上重标定 $\tau$， scorer 代码保持与主仓一致。
光照差异会导致 $\tau$ 漂移；报告应同时给出本地与 ShezhenV3 全局中位数。

\paragraph{安全与隐私.}
图像在 \texttt{upload\_cloud} 之前留于边缘；API Key 保护运动与视觉接口。
去标识 JSONL 不含患者元数据；IQA 路径仅使用几何框，不使用诊断标签。

\paragraph{与 MoveIt 技能集成.}
具名位姿（\texttt{tongue\_pose}、\texttt{face\_pose}）解耦采集几何与 IQA；可在最终舌采前于 \texttt{face\_pose} 做预检。

\paragraph{回归测试.}
单元测试覆盖 flags 推导；全矩阵回归见 \texttt{paper1\_run\_tcm\_full.sh}（含 B3 训练时约 10--60 分钟）。

\paragraph{文档地图.}
英文 \texttt{docs/}、中文学习手册 \texttt{docs-zh/paper1/study/}（六章）与 \texttt{REPRODUCE.md} 构成新人上手路径。
